\section{Finite Automata}

\subsection{Symbol, Alphabets and Words}

\begin{definition}
    \textbf{Alphabet:} finite, non-empty set.
\end{definition}

\begin{example}
    \textbf{Examples of Alphabets:}
    \begin{itemize}
        \item Boolean alphabet ${0,1}$
        \item latin alphabet ${a,b,\dots,z}$
        \item logical alphabet ${0,a,x,(,),\dots}$
    \end{itemize}
\end{example}

\begin{definition}
    \textbf{Word:} A word over an alphabet $A$ is a finite sequence of symbols
    of $A$.
\end{definition}

\begin{definition}
    The \textbf{lenght} of a word it the number of symbols it uses. There is
    always an empty word of length zero.
\end{definition}

\begin{remark}
    Words can be concatenated or rewersed to form new words.
    $$
        kon(x,y) = xy
    $$
    $$
        w = a_1 a_2 a_3 \implies w^R = a_3 a_2 a_1
    $$
\end{remark}

\begin{definition}
    The $i$-th iteration of a word $x = ab$ is $x^i = ababab \dots ab$ the word
    concatenated with itself i times.
\end{definition}

\begin{definition}
    A word $v$ is a partial-, prefix- or suffix-word of a word $w$ if $v$ is
    contained in $w$, $w$ begins with $v$ or ends with $v$.
\end{definition}

\begin{remark}
    For a word $x$ and symbol $a$ the expression $|x|_a$ describes the numnber
    of apperances of $a$ in $x$.
\end{remark}

\begin{theorem}
    For an alphabet with an order imposed on its symbols ($a_1 \le a_2 \le ...$)
    we can define the harmonic order of its words as a lexicographic order where
    $$
        u \le  v
          \iff (|u| \le |v|)
               \lor (|u| = |v| \land u = xs_iu',v = xs_jv'. s_i \le s_j)
    $$
\end{theorem}

\subsection{Formal Languages}

\begin{definition}
    A subset $L \subset \Sigma^*$ is called Language over the alphabet $\Sigma$
\end{definition}

\begin{definition}
    The \textbf{complement} $L^C$ of a Language $L$ is definied as
    $\Sigma^* \setminus L$
\end{definition}

\begin{definition}
    The \textbf{concatenation} of two languages $L_1, L_2$ is all the set of all
    $w_1w_2$ where $w_1 \in L_1, w_2 \in L_2$
\end{definition}

\begin{remark}
    We can concatenate a Language with itself (possibly multiple times or zero)
    $$
        L^0 = \{\lambda\}, L^1
            = L, L^2
            = LL \dots L^*
            = \bigcup_{i=0}^\infty L^i, L^+
            = LL^*
    $$
\end{remark}

\begin{lemma}
    Let $L_1,L_2, L_3$ be languages over an alphabet $\Sigma$ then
    $$
        L_1L_2 \cup L_1L_3 = L_1(L_2 \cup L_3)
    $$
    $$
        L_1(L_2 \cap L_3) \subseteq L_1L_2 \cap L_1L_3
    $$
\end{lemma}

\subsection{Algorithmic Problems}

\begin{remark}
    An algorithm describes a function $\Sigma_{in}^* \to \Sigma_{out}^*$
\end{remark}

\begin{definition}
    Let $\Sigma$ be an alphabet and $L \subseteq \Sigma^*$. The decision problem
    $(\Sigma,L)$ is defined (for arb $x \in Sigma^*$) as
    $$
        x \in L \quad \text{or} \quad x \notin L
    $$
\end{definition}

\begin{remark}
    If an algorithm A correctly decides for any $x$ then we say
    $$
        A \quad \text{detects the Language} \quad L
    $$
\end{remark}

\subsection{Finite Automata}

\begin{remark}
    A finite automata (finite state machine) is a set of states together with
    state transitions on inputs with fixed boolean outputs on each state.
\end{remark}

\begin{definition}
    \textbf{Finite Automata:} is a Quadroupel $M = (Q,\Sigma, \delta, q_0, F)$
    \begin{itemize}
        \item $Q$ is a finite set of states
        \item $\Sigma$ is an alphabet (Input-alphabet)
        \item $\delta$ is a function $Q \times \Sigma \to Q$
        \item $q_0 \in Q$ is the starting-state
        \item $F \subseteq Q$ is the set of accepting (true) states
    \end{itemize}
\end{definition}

\begin{definition}
    \textbf{Configuration:} of an Automata is an Elment of $Q \times \Sigma^*$
    \begin{itemize}
        \item \textbf{Starting-Configuraiton:} $C_0 = (q_0, x) | x \in \Sigma^*$
        \item \textbf{End-Configruation:} $C_n = (q_0, \lambda)$ where its
              accepting if $C_n \in F \times \{\lambda\}$
    \end{itemize}
\end{definition}

\begin{remark}
    A calculation can be seen as a finite sequence of configurations
    $C_1, C_2 \dots$ such that for every $C_i \to_{s_M} C_{i+1}$ with steps
    $s_M$ in our machine $M$
    $$
        (q_0, 01001) \vdash_{M} (q_x, 1001) \vdash_{M} (q_y, 001) \dots
    $$
    where a step is a relation on configurations
    $$
        \vdash_{M} \subseteq (Q \times \Sigma^*) \times (Q \times \Sigma^*)
    $$
\end{remark}

\begin{definition}
    The \textbf{transitive hull of transitions} describes the calculation over
    multiple steps:
    $$
        (q,w) \vdash_{M}^{*} (p,u)
              \iff           (q = p \land w = u) \lor \exists k
              \in            \mathbb{N}^{+} \text{ s.t. }
    $$
    $$
        w =   a_1 \dots a_k u \quad \text{with} \quad a_i
          \in \Sigma \quad \text{and}
    $$
    $$
        \exists r_1, \dots, r_{k-1} \in        Q \text{ s.t. } (q,w)
                                    \vdash_{M} (r_1, a_2 \dots a_k u)
                                    \vdash_{M} (r_2, a_3 \dots a_k u) \dots
                                    \vdash_{M} (p,u)
    $$
\end{definition}

\begin{definition}
    The \textbf{extended transition function} evaluates the end state of a word
    over multiple transitions:
    $$
        \tilde{\delta}: Q \times \Sigma^{*} \to Q
    $$
    $$
        \tilde{\delta}(q, \lambda) = q \quad \forall q \in Q
    $$
    $$
        \tilde{\delta}(q, wa)
        =   \delta(\tilde{\delta} (q,w), a) \quad \forall a
        \in \Sigma, w
        \in \Sigma^{*}, q
        \in Q
    $$
\end{definition}

\begin{remark}
    States can be grouped into equivalence classes containing all elements with
    the same end state. This is useful for proofs regarding the language of
    finite automata, which are typically done through these classes and
    induction.
\end{remark}

\begin{definition}
    \textbf{Product Automata:} Used to determine the intersection, union or
    exclusion ($\odot$) of two finite automata (their languages), built with the
    state space $Q = Q_1 \times Q_2$.
    $$
        M_1 = (Q_1, \Sigma, \delta_1, q_{01}, F_1), \quad M_2
            = (Q_2, \Sigma, \delta_2, q_{02}, F_2)
    $$
    $$
        \implies M
        =        (Q_1 \times Q_2, \delta((x,y), a) = (\delta_1(x, a), \delta_2(y, a)), (q_01, q_02), F_1 \odot F_2)
    $$
\end{definition}

\begin{remark}
    The correctness of this product automata can be proven by induction over the
    word length.
\end{remark}

\subsection{Non-regularity}

\begin{definition}
    A language $L$ is called \textbf{regular} if there exists a finite automata
    that accepts/detects the language $L$ and \textbf{non-reular} else.
\end{definition}

\begin{lemma}
    Let $A = (Q, \Sigma, \delta, q_0, F)$ be a finite automata and let
    $x,y \in \Sigma^*$ such that
    $$
        \tilde{\delta}(q_0,x) = \tilde{\delta}(q_0, y)
    $$
    then there exists $z \in \Sigma^*$ such that
    $$
        Kl[xz] = Kl[yz]
    $$
\end{lemma}

\begin{remark}
    This illustrates that finite machines are finite hence can not distinguish
    beetween $xz$ and $yz$ once they are in a certain state $q_r$, they do not
    know how they got there.
\end{remark}

\begin{proof}
    \textbf{Proof of Non-Regularity:} Let $L = \{0^n1^n | n \in \mathbb{N}\}$
    and assume $L \in L_{EA}$ ($L$ is regular) hence there exists a finite
    automata $A$ such that $L(A) = L$. Let $m = |Q|$ and observe the words
    $x \in \{0^n | n \leq m + 1 \}$. Because this are more words than $A$ has
    states there exists $i \neq j \in [0, m + 1] \cap \mathbb{N}$ such that
    $$
        \tilde{\delta}(q_0, 0^i) = \tilde{\delta}(q_0, 0^j)
    $$
    Hence by the above lemma the words $0^i1^i, 0^j1^i$ will both be accepted
    which contradicts the assumption that $A$ detects $L$ because
    $0^j1^i \notin L$.
\end{proof}

\begin{lemma}
    \textbf{Pumping-Lemma:} Let a language $L \in \Sigma^*$ be regular, then
    there exists a constant $n_0 \in \mathbb{N}$ such that every word
    $w \in \Sigma^*$ with $|w| \ge n_0$ can be taken apart into $x,y,z$ where
    \begin{itemize}
        \item $|yx| \le n_0$
        \item $|x| \ge 1$
        \item $\{yx^kz | k \in \mathbb{N} \} \subseteq L \lor \{yx^kz | k \in \mathbb{N} \} \cap L = \emptyset$
    \end{itemize}
\end{lemma}

\begin{example}
    $$
        L = {0^n 1^n | n \in \mathbb{N}}
    $$

    \textbf{Proof using Pumping Lemma:} Assume $L$ regular hence there exists
    constant $n_0$ with the properties described by the pumping lemma.
    $$
        w = 0^{n_0} 1^{n_0}
    $$
    And there exists $yxz = w$ that fulfill the properties of the pumping lemma.
    Hence
    $$
        |yx| \le n_0 \quad \implies y = 0^e, x = 0^m \quad e,m \in \mathbb{N}
    $$
    $$
        |x| \ge 1 \quad \implies m \ge 1
    $$
    $$
        |yxz| =         n_0
              \implies  \{0^{n_0-m+km} 1^n_0 | k \in \mathbb{N} \}
              \subseteq L
    $$
    Observe that this is a contradiction as for $k = 0$ the last condition does
    not hold.
    $$
        yx^0z = yz = 0^{n_0-m}1^{n_0} \notin L
    $$
\end{example}

\subsection{Nondeterminism}

\begin{important}
    Nondeterminism has nothing to do with Non Deterministic (random) behavior.
\end{important}

\begin{definition}
    Nondeterministic finite automata are finite automatas where the
    transition-function $\delta$ is not well definied hence
    $$
        \hat{\delta}(q_k, w) = \{q_i, q_j ...\}
    $$
    and a
    $$
        w \in L \iff \exists q_l \in F \cap \{q_i, q_j ...\}
    $$
    Hence a word is accepted if there is at least one valid path that leads to
    an accepting endstate
\end{definition}

\subsection{Power-set-construction}

\begin{theorem}
    \textbf{Power-Set Construction (Subset Construction):} For every
    nondeterministic finite automaton $M = (Q_N, \Sigma, \delta_N, q_0, F_N)$,
    there exists an equivalent deterministic finite automaton
    $M_D = (Q_D, \Sigma, \delta_D, q_{0,D}, F_D)$ such that
    $$
        L(M_D) = L(M)
    $$
\end{theorem}

\begin{definition}
    Given an NFA $M = (Q_N, \Sigma, \delta_N, q_0, F_N)$, the equivalent DFA
    $M_D = (Q_D, \Sigma, \delta_D, q_{0,D}, F_D)$ is constructed as follows:
    \begin{itemize}
        \item \textbf{States:} $Q_D = \mathcal{P}(Q_N)$ (the power set of $Q_N$,
              i.e., every state in $M_D$ is a subset $S \subseteq Q_N$)
        \item \textbf{Alphabet:} $\Sigma$ (identical to the NFA's alphabet)
        \item \textbf{Start State:} $q_{0,D} = \{q_0\}$
        \item \textbf{Transition Function:} For every $S \in Q_D$ and
              $a \in \Sigma$:
              $$
                  \delta_D(S, a) = \bigcup_{q \in S} \delta_N(q, a)
              $$
        \item \textbf{Accepting States:} A state $S \in Q_D$ is accepting if it
              contains at least one accepting state of $M$:
              $$
                  F_D = \{S \subseteq Q_N \mid S \cap F_N \neq \emptyset\}
              $$
    \end{itemize}
\end{definition}

\begin{remark}
    The intuition behind the construction is that the DFA tracks all possible
    states the NFA could currently be in simultaneously.
    While the power set has $|Q_D| = 2^{|Q_N|}$ states, in practice only the
    states reachable from $\{q_0\}$ need to be constructed.
    By induction on word length $|w|$, one can show that
    $$
        \tilde{\delta}_D(\{q_0\}, w) = \hat{\delta}_N(q_0, w)
    $$
    which immediately implies $w \in L(M_D) \iff w \in L(M)$.
\end{remark}

\begin{theorem}
    \textbf{State Explosion / Exponential Lower Bound:} For every $k \in \mathbb{N}^+$,
    let $L_k$ be the language over $\Sigma = \{0, 1\}$ defined by:
    $$
        L_k = \{ w \in \{0, 1\}^* \mid \text{the } k\text{-th symbol from the end of } w \text{ is } 1 \}
    $$
    Then:
    \begin{enumerate}
        \item There exists an NFA accepting $L_k$ with only $k + 1$ states.
        \item Every DFA accepting $L_k$ requires at least $2^k$ states.
    \end{enumerate}
\end{theorem}

\begin{example}
    \textbf{NFA for $L_k$:} We can construct an NFA $M = (Q, \{0, 1\}, \delta, q_0, \{q_k\})$
    with $Q = \{q_0, q_1, \dots, q_k\}$:
    \begin{itemize}
        \item $\delta(q_0, 0) = \{q_0\}$
        \item $\delta(q_0, 1) = \{q_0, q_1\}$ (nondeterministically guess that this $1$ is the $k$-th symbol from the end)
        \item $\delta(q_i, 0) = \delta(q_i, 1) = \{q_{i+1}\}$ for $1 \le i \le k - 1$
    \end{itemize}
    This automaton has exactly $k + 1$ states and accepts $L_k$.
\end{example}

\begin{proof}
    \textbf{Proof ntuition for the $2^k$ DFA Lower Bound:}
    Unlike an NFA, a deterministic automaton cannot guess when it is $k$ symbols away from the end. 
    Since the input could stop at any time, a DFA must maintain the exact history of the last $k$ symbols seen. There are $2^k$ possible words of length $k$, so it must distinguish all of them.
\end{proof}
